%% ----------------------------------------------------------------------------
%%  §3.1  Premessa: l'area del cerchio e l'idea di limite
%% ----------------------------------------------------------------------------
\section{Premessa: l'area del cerchio}
\label{sec:premessa-limiti}

Uno dei problemi che affronteremo nel capitolo sugli integrali sarà il calcolo delle
aree di figure piane. Tale calcolo fa uso in modo essenziale del concetto di limite. Per
mostrarlo, consideriamo un metodo per il calcolo dell'area di un cerchio simile a quello
che Archimede usò nel III secolo a.C.

Si tratta di «approssimare» un cerchio di raggio $1$ (potremmo scegliere anche un
diverso valore del raggio) con figure che «differiscono di poco» dal cerchio e di cui sia
possibile calcolare l'area. Nel capitolo sugli integrali vedremo che è comodo
considerare, come figure approssimanti, unioni di rettangoli. Qui, per ogni $n\ge3$,
consideriamo il poligono regolare $P_{n}$ di $n$ lati inscritto nel cerchio di raggio
$1$: per $n=3$ un triangolo equilatero, per $n=4$ un quadrato, per $n=6$ un esagono, e
così via (\cref{fig:poligoni-inscritti}).

\begin{figure}[htb]
\centering
\begin{tikzpicture}[scale=1.0, every node/.style={font=\footnotesize}]
  \foreach \n/\xx/\col in {3/0/blu, 4/2.45/rame, 6/4.9/verde} {
    \begin{scope}[shift={(\xx,0)}]
      \draw[grigio!70] (0,0) circle (1);
      \draw[\col, thick, fill=\col!12] (90:1) \foreach \k in {1,...,\n} { -- ({90+360*\k/\n}:1)} -- cycle;
      \fill (0,0) circle (0.8pt);
      \node at (0,-1.3) {$P_{\n}$};
    \end{scope}}
  % dettaglio di un triangolo
  \begin{scope}[shift={(6.6,0)}, scale=1.75]
    \pgfmathsetmacro{\al}{30}
    \draw[grigio!70] ({-\al-10}:1) arc ({-\al-10}:{\al+10}:1);
    \draw[blu, thick] (0,0) -- (\al:1) -- (-\al:1) -- cycle;
    \draw[blu, dashed] (0,0) -- ({cos(\al)},0);
    \draw[blu] ({cos(\al)},0) -- (\al:1);
    \draw ({cos(\al)-0.06},0) -- ({cos(\al)-0.06},0.06) -- ({cos(\al)},0.06);
    \draw[rame] (0.3,0) arc (0:\al:0.3); \node[rame] at (15:0.45) {$\alpha$};
    \node[below] at ({cos(\al)/2},-0.02) {$\cos\alpha$};
    \node[right] at ({cos(\al)},{sin(\al)/2}) {$\sen\alpha$};
    \node[above left] at ({cos(\al)/2},{sin(\al)/2}) {$1$};
    \node at (0.55,-0.85) {$\alpha=\pi/n$};
  \end{scope}
\end{tikzpicture}
\caption{I poligoni regolari $P_{3}$, $P_{4}$, $P_{6}$ inscritti nel cerchio di raggio
$1$. A destra, uno dei $2n$ triangoli rettangoli in cui si scompone $P_{n}$: l'angolo al
centro vale $\alpha=\pi/n$, i cateti valgono $\sen\alpha$ e $\cos\alpha$.}
\label{fig:poligoni-inscritti}
\end{figure}

Agli $n$ lati del poligono $P_{n}$ corrispondono $n$ angoli al centro uguali fra loro,
che sono l'$n$-esima parte dell'angolo giro e che quindi misurano $2\pi/n$ radianti.
Tracciando le altezze dal centro, $P_{n}$ si scompone in $2n$ triangoli rettangoli
uguali, con l'angolo al centro $\alpha=\pi/n$ e ipotenusa $1$; i cateti misurano quindi
$\sen\alpha$ e $\cos\alpha$, e l'area di ciascun triangolo vale
$\frac12\sen\alpha\cos\alpha$. Indicando con $a_{n}$ l'area di $P_{n}$ otteniamo
\begin{equation}\label{eq:area-poligono-1}
  a_{n}=\operatorname{area}(P_{n})=2n\cdot\frac12\sen\alpha\cos\alpha,
\end{equation}
e, usando la formula di duplicazione $\sen2\alpha=2\sen\alpha\cos\alpha$
\eqref{eq:duplicazione} e ricordando che $\alpha=\pi/n$,
\begin{equation}\label{eq:area-poligono}
  a_{n}=\frac n2\,\sen\frac{2\pi}{n}.
\end{equation}
Riportiamo $a_{n}$ per alcuni valori di $n$ nella \cref{tab:aree-poligoni}; il lettore
controlli i primi valori, fino a $n=12$, usando la \cref{tab:valori-notevoli} (per $n=3$
si ricordi che $\sen\frac{2\pi}3=\sen\bigl(\pi-\frac\pi3\bigr)=\sen\frac\pi3$).

\begin{table}[htb]
\centering
\renewcommand{\arraystretch}{1.5}
\footnotesize\setlength{\tabcolsep}{4pt}
\begin{tabular}{l*{10}{c}}
\toprule
$n$ & $3$ & $4$ & $6$ & $8$ & $12$ & $20$ & $50$ & $100$ & $200$ & $1000$\\
\midrule
$a_{n}$ & $\tfrac{3\sqrt3}{4}$ & $2$ & $\tfrac{3\sqrt3}{2}$ & $2\sqrt2$ & $3$ & $3{,}0902$ &
$3{,}1333$ & $3{,}1395$ & $3{,}1411$ & $3{,}14157$\\
\bottomrule
\end{tabular}
\caption{Le aree $a_{n}$ dei poligoni regolari di $n$ lati inscritti nel cerchio di
raggio $1$, calcolate con la \eqref{eq:area-poligono}.}
\label{tab:aree-poligoni}
\end{table}

Come ci aspettiamo, per $n$ grande i valori di $a_{n}$ si avvicinano a
$\pi=3{,}14159265\ldots$ Alla fine del \cref{sec:limiti-notevoli} dimostreremo, tramite
la \eqref{eq:area-poligono}, che l'area del cerchio di raggio $1$ vale $\pi$. La
\eqref{eq:area-poligono} non è però utile per il calcolo \emph{numerico} di $\pi$, perché
$\pi$ entra nella definizione stessa di $a_{n}$ (Archimede, come vedremo nella nota
storica, aggirò il problema in modo geniale).

Qui desideriamo sottolineare lo schema del procedimento: abbiamo introdotto una
\emph{successione} di numeri reali
\begin{equation}\label{eq:succ-aree}
  a_{3},\ a_{4},\ a_{5},\ a_{6},\ \dots,\ a_{n},\ \dots
\end{equation}
Di tale successione ci interessa il comportamento per $n$ grande. Infatti il numero
$a_{n}$ rappresenta l'area del poligono regolare di $n$ lati inscritto nel cerchio, ed
è tanto più vicino all'obiettivo --- l'area del cerchio --- quanto più $n$ è grande. Se
avesse un senso, diremmo che ci interessa «l'ultimo» termine della successione; però ciò
non ha senso, perché non esiste l'ultimo termine della successione.

Ciò che realmente ci interessa è il \emph{limite} della successione $a_{n}$, cioè un
numero $a\in\R$ che sia «vicino» ai termini della successione che hanno l'indice $n$
«grande». Ciò si esprime più precisamente così: $a$, limite della successione $a_{n}$,
è un numero reale tale che, comunque si scelga un intervallo di numeri intorno ad $a$,
diciamo $(a-\epsilon,a+\epsilon)$ con $\epsilon>0$, esiste un indice $\nu$ tale che, per
$n>\nu$, $a_{n}$ sta in questo intervallo, cioè $a-\epsilon<a_{n}<a+\epsilon$.

\begin{figure}[htb]
\centering
\begin{tikzpicture}[x=12cm, every node/.style={font=\footnotesize}]
  % a_n = 1 - 0.9 (-1)^n / n ;  |a_n - 1| = 0.9/n < eps  sse  n > 0.9/eps
  \foreach \row/\eps/\nuu/\lab in {0/0.30/3/{0{,}30}, 1/0.12/7/{0{,}12}, 2/0.04/22/{0{,}04}} {
    \begin{scope}[yshift={-\row*1.25cm}]
      \fill[verde!18] ({1-\eps},-0.16) rectangle ({1+\eps},0.16);
      \draw[-Latex, grigio] (0.70,0) -- (1.40,0);
      \foreach \n in {3,...,60} {
        \pgfmathsetmacro{\x}{1-0.9/\n*cos(180*\n)}
        \pgfmathparse{\n>\nuu ? 1 : 0}
        \ifnum\pgfmathresult=1 \fill[verde!70!black] (\x,0) circle (1.1pt);
        \else \fill[rame] (\x,0) circle (1.4pt); \fi
      }
      \draw[verde!60!black, thick] ({1-\eps},-0.16) -- ({1-\eps},0.16);
      \draw[verde!60!black, thick] ({1+\eps},-0.16) -- ({1+\eps},0.16);
      \node[verde!60!black, above] at ({1-\eps},0.16) {$a-\epsilon$};
      \node[verde!60!black, above] at ({1+\eps},0.16) {$a+\epsilon$};
      \draw[blu, thick] (1,-0.1) -- (1,0.1); \node[blu, below=2pt] at (1,-0.1) {$a$};
      \node[anchor=west] at (1.41,0) {$\epsilon=\lab$, \ $\nu=\nuu$};
    \end{scope}}
\end{tikzpicture}
\caption{L'idea di limite, per la successione $a_{n}=1-\frac{9}{10}\frac{(-1)^{n}}{n}$ e $a=1$
(qui $\abs{a_{n}-a}<\epsilon$ se e solo se $n>\frac{9}{10\epsilon}$). Per tre valori sempre
più piccoli di $\epsilon$, i termini
$a_{n}$ con $n>\nu$ (in verde) cadono tutti nell'intervallo $(a-\epsilon,a+\epsilon)$;
quelli con $n\le\nu$ (in rame) possono cadere fuori. Diminuendo $\epsilon$, l'indice
$\nu$ in genere aumenta; l'importante è che, \emph{comunque} si fissi $\epsilon>0$, un
tale $\nu$ esista.}
\label{fig:idea-limite}
\end{figure}

Abbiamo rappresentato nella \cref{fig:idea-limite} tre casi per diversi valori di
$\epsilon$; se si cambia $\epsilon$ con un numero minore, l'indice $\nu$ in genere
aumenta. L'importante è che, comunque si fissi $\epsilon>0$, in corrispondenza esista un
indice $\nu$ con le proprietà suddette.

\begin{definizione}[Intorno]\index{intorno}
Dati $a\in\R$ ed $\epsilon>0$, l'intervallo
\begin{equation}\label{eq:intorno}
  (a-\epsilon,a+\epsilon)=\{x\in\R : a-\epsilon<x<a+\epsilon\}=\{x\in\R : \abs{x-a}<\epsilon\}
\end{equation}
si chiama \emph{intorno} del numero $a$ di raggio $\epsilon$. L'ultima uguaglianza è la
\eqref{eq:mod-lt}: l'intorno di raggio $\epsilon$ è l'insieme dei punti che distano da
$a$ meno di $\epsilon$.
\end{definizione}

\begin{notastorica}[Archimede e il calcolo di $\pi$]
Nella \emph{Misura del cerchio} Archimede (287--212 a.C.) dimostrò che l'area del
cerchio è uguale a quella di un triangolo rettangolo con un cateto uguale al raggio e
l'altro uguale alla circonferenza --- in linguaggio moderno, $\pi r^{2}=\frac12\,r\cdot2\pi
r$ --- con un raffinato doppio ragionamento per assurdo (il \emph{metodo di esaustione})
che è, in sostanza, un passaggio al limite. Nella stessa opera calcolò il perimetro dei
poligoni regolari inscritti e circoscritti partendo dall'esagono e raddoppiando quattro
volte il numero dei lati, fino a $96$; ottenne così
\[
  3+\tfrac{10}{71}<\pi<3+\tfrac17,\qquad\text{cioè}\qquad 3{,}1408<\pi<3{,}1429 .
\]
Il punto geniale è che il procedimento \emph{non} richiede di conoscere $\pi$: le
grandezze dei poligoni con $2n$ lati si ottengono da quelle dei poligoni con $n$ lati
mediante sole operazioni aritmetiche ed estrazioni di radice quadrata (si veda il codice
qui sotto). Archimede eseguì a mano tutte le radici con approssimazioni razionali,
sempre per difetto o per eccesso nel verso giusto, per non compromettere le
disuguaglianze.
\end{notastorica}

\begin{python}[Il metodo di Archimede, senza conoscere $\pi$]
# s = semiperimetro del poligono inscritto, t = del circoscritto
# (raggio 1). Da n a 2n lati (media armonica e media geometrica):
#     t_2n = 2 s t / (s + t),      s_2n = sqrt(s * t_2n)
# Si parte dall'esagono: s_6 = 3, t_6 = 2*sqrt(3). Niente pi greco!
import math

n, s, t = 6, 3.0, 2 * math.sqrt(3)
print(f"{n:4d} lati: {s:.6f} < pi < {t:.6f}")
for _ in range(4):               # quattro raddoppi, fino a 96 lati
    t = 2 * s * t / (s + t)
    s = math.sqrt(s * t)
    n *= 2
    print(f"{n:4d} lati: {s:.6f} < pi < {t:.6f}")
print(f"Archimede: {223/71:.6f} < pi < {22/7:.6f}")
\end{python}
\begin{risultato}
   6 lati: 3.000000 < pi < 3.464102
  12 lati: 3.105829 < pi < 3.215390
  24 lati: 3.132629 < pi < 3.159660
  48 lati: 3.139350 < pi < 3.146086
  96 lati: 3.141032 < pi < 3.142715
Archimede: 3.140845 < pi < 3.142857
\end{risultato}

Le due successioni del programma, quella dei semiperimetri inscritti (crescente) e
quella dei circoscritti (decrescente), si stringono l'una contro l'altra intorno a $\pi$:
è una situazione che ritroveremo, in forma generale, nel teorema sulle successioni
monotòne (\cref{sec:monotone}) e nella definizione del numero $e$ (\cref{sec:numero-e}).
